\subsection*{CU-33: Silenciar o bloquear a un jugador}
\addcontentsline{toc}{subsection}{CU-33: Silenciar o bloquear a un jugador}
\label{cu-33}

\begin{fichacu}{CU-33}{Silenciar o bloquear a un jugador}
  \crudFila{Responsable}{Ángel Gabriel Aguilar Hernández}
  \crudFila{Actor principal}{Jugador}
  \crudFila{Actores de sistema}{Servidor de partidas, Base de datos}
  \crudFila{Prototipos}{10a, 10c}
  \crudFila{Objetivo}{Permitir al Jugador dejar de leer a otro jugador —silenciándolo— o cortar además todo contacto con él —bloqueándolo—, y deshacer cualquiera de las dos cosas desde sus ajustes.}
  \crudFila{Disparador}{El Jugador selecciona ``silenciar'' o ``bloquear'' en el menú de otro jugador, o abre la lista de silenciados y bloqueados en la \texttt{GUI\_\allowbreak{}Settings}.}
  \textbf{Precondiciones} & \cuCelda{\begin{cureglas}
      \item \textbf{\mbox{PRE-01}} El Jugador tiene una \textit{Sesion} abierta y su token es válido.
      \item \textbf{\mbox{PRE-02}} El \textit{Usuario} tiene \texttt{tipo\_\allowbreak{}cuenta} en \texttt{REGISTRADA}. \textit{(Ver \mbox{FA-01})}
      \item \textbf{\mbox{PRE-03}} El Servidor de partidas está en ejecución y tiene conexión disponible con la Base de datos.
    \end{cureglas}} \\ \hline
  \textbf{Flujo normal} & \cuCelda{\begin{cuenum}[start=1]
      \item El Jugador abre el menú de otro jugador desde el chat, el ranking, la lista de amigos o la partida, y selecciona ``bloquear''. \textit{(Ver \mbox{FA-02})}
      \item El SJM despliega la \texttt{GUI\_\allowbreak{}MessageConfirm} explicando qué implica: no podrán emparejarlos, invitarse, buscarse ni escribirse, y la amistad se eliminará si existe; con las opciones ``Bloquear'' y ``Cancelar''. \textit{(Ver \mbox{FA-03})}
      \item El Jugador selecciona ``Bloquear''.
      \item El SJM envía el mensaje \texttt{BLOCK\_\allowbreak{}PLAYER\_\allowbreak{}REQUEST} con el \texttt{id\_\allowbreak{}usuario} del otro jugador y el token. \textit{(Ver \mbox{EX-02}, \mbox{EX-03}, \mbox{EX-04})}
      \item El Servidor de partidas verifica que el token corresponda a una \textit{Sesion} abierta y que el otro jugador exista y no sea el propio Jugador. \textit{(Ver \mbox{FA-04})}
      \item El Servidor de partidas verifica que no exista ya un \textit{Bloqueo} del Jugador hacia ese jugador. \textit{(Ver \mbox{FA-05})}
      \item El Servidor de partidas inicia una transacción y crea el \textit{Bloqueo} con el \texttt{id\_\allowbreak{}bloqueador}, el \texttt{id\_\allowbreak{}bloqueado} y la \texttt{fecha\_\allowbreak{}bloqueo}. \textit{(Ver \mbox{EX-01})}
      \item En la misma transacción, el Servidor de partidas elimina la \textit{Amistad} entre ambos si existe, elimina las \textit{Solicitud} pendientes en cualquier sentido y marca como \texttt{EXPIRADA} cualquier \textit{Invitacion} pendiente entre ellos (\mbox{RN-03}).
    \end{cuenum}} \\
   & \cuCelda{\begin{cuenum}[start=9]
      \item Si ambos están en la \textit{ColaEmparejamiento}, el Servidor de partidas impide que se emparejen a partir de este momento; si ambos están en la misma \textit{Sala} abierta, retira de ella al bloqueado cuando el anfitrión es el Jugador, o retira al Jugador en caso contrario (\mbox{RN-05}).
      \item El Servidor de partidas confirma la transacción y \textbf{no} notifica nada al bloqueado (\mbox{RN-02}). \textit{(Ver \mbox{EX-01})}
      \item El Servidor de partidas responde con \texttt{BLOCK\_\allowbreak{}PLAYER\_\allowbreak{}OK}.
      \item El SJM oculta los mensajes del bloqueado en todos los canales, lo retira de la lista de amigos si estaba y muestra una confirmación breve. \textit{(Ver \mbox{FA-06})}
      \item Fin del caso de uso.
    \end{cuenum}} \\ \hline
  \textbf{Flujos alternos} & \cuCelda{\textbf{\mbox{FA-01}: El Jugador es un invitado}\par
    \begin{cuenum}[start=1]
      \item El SJM ofrece ``silenciar'' pero no ``bloquear'', porque el bloqueo exige una cuenta rastreable, y junto a la opción deshabilitada propone vincular la cuenta.
      \item Si el Jugador silencia, el flujo continúa en el \mbox{FA-02}.
    \end{cuenum}} \\
   & \cuCelda{\textbf{\mbox{FA-02}: El Jugador silencia en lugar de bloquear}\par
    \begin{cuenum}[start=1]
      \item El Jugador selecciona ``silenciar''.
      \item El SJM envía el mensaje \texttt{MUTE\_\allowbreak{}PLAYER\_\allowbreak{}REQUEST} sin pedir confirmación, porque silenciar es inmediato y reversible (\mbox{RN-01}).
      \item El Servidor de partidas crea el \textit{Silencio} con el \texttt{id\_\allowbreak{}usuario}, el \texttt{id\_\allowbreak{}silenciado} y la \texttt{fecha\_\allowbreak{}silencio}, si no existía.
      \item El Servidor de partidas responde con \texttt{MUTE\_\allowbreak{}PLAYER\_\allowbreak{}OK}.
      \item El SJM deja de mostrar los mensajes y los emotes del silenciado en todos los canales, y muestra la opción ``dejar de silenciar'' en su menú.
      \item Fin del caso de uso. Pueden seguir coincidiendo en partidas.
    \end{cuenum}} \\
   & \cuCelda{\textbf{\mbox{FA-03}: El Jugador cancela el bloqueo}\par
    \begin{cuenum}[start=1]
      \item El Jugador selecciona ``Cancelar''.
      \item El SJM cierra la confirmación sin enviar nada.
      \item Fin del caso de uso.
    \end{cuenum}} \\
   & \cuCelda{\textbf{\mbox{FA-04}: El jugador indicado no existe o es el propio Jugador}\par
    \begin{cuenum}[start=1]
      \item El Servidor de partidas responde con \texttt{BLOCK\_\allowbreak{}PLAYER\_\allowbreak{}FAILED} y el motivo \texttt{JUGADOR\_\allowbreak{}INVALIDO}.
      \item Si el jugador indicado es el propio, el Servidor de partidas lo anota en la bitácora del servidor como indicio de cliente modificado.
      \item Fin del caso de uso.
    \end{cuenum}} \\
   & \cuCelda{\textbf{\mbox{FA-05}: El jugador ya estaba bloqueado}\par
    \begin{cuenum}[start=1]
      \item El Servidor de partidas responde con \texttt{BLOCK\_\allowbreak{}PLAYER\_\allowbreak{}OK} sin modificar nada, porque el estado pedido ya se cumple.
      \item El flujo continúa en el paso 12 del FN.
    \end{cuenum}} \\
   & \cuCelda{\textbf{\mbox{FA-06}: El Jugador bloquea a su rival durante una partida}\par
    \begin{cuenum}[start=1]
      \item El Servidor de partidas aplica el bloqueo, pero la partida en curso \textbf{sigue} hasta terminar (\mbox{RN-04}).
      \item El SJM del Jugador oculta desde ese momento el chat y los emotes del rival en esa partida.
      \item Al terminar, no podrán volver a emparejarse.
      \item Fin del caso de uso.
    \end{cuenum}} \\
   & \cuCelda{\textbf{\mbox{FA-07}: El Jugador deshace un silencio o un bloqueo}\par
    \begin{cuenum}[start=1]
      \item El Jugador abre la lista de silenciados y bloqueados en la \texttt{GUI\_\allowbreak{}Settings}.
      \item El SJM envía \texttt{LIST\_\allowbreak{}MUTES\_\allowbreak{}BLOCKS\_\allowbreak{}REQUEST} y el Servidor de partidas responde con los \textit{Silencio} y \textit{Bloqueo} del Jugador, con el \texttt{nickname} y la fecha de cada uno.
      \item El Jugador selecciona ``dejar de silenciar'' o ``desbloquear'' en uno de ellos.
      \item El Servidor de partidas elimina la fila de \textit{Silencio} o de \textit{Bloqueo} correspondiente.
      \item El SJM lo retira de la lista. La amistad eliminada al bloquear \textbf{no} se restaura: habría que volver a solicitarla (\mbox{RN-06}).
      \item Fin del caso de uso.
    \end{cuenum}} \\
   & \cuCelda{\textbf{\mbox{FA-08}: El Jugador bloquea a un moderador o es bloqueado por un reportado}\par
    \begin{cuenum}[start=1]
      \item El Servidor de partidas aplica el bloqueo con normalidad.
      \item El bloqueo no afecta a las funciones de moderación: un moderador sigue pudiendo revisar reportes y sancionar a quien lo bloqueó, y un reportado no puede impedir que se revise su conducta bloqueando a su denunciante (\mbox{RN-07}).
      \item Fin del caso de uso.
    \end{cuenum}} \\ \hline
  \textbf{Excepciones} & \cuCelda{\textbf{\mbox{EX-01}: Fallo al escribir en la base de datos}\par
    \begin{cuenum}[start=1]
      \item El Servidor de partidas revierte la transacción, de modo que no quede el \textit{Bloqueo} creado con la \textit{Amistad} todavía vigente.
      \item El Servidor de partidas registra la excepción con un identificador de incidencia y responde con \texttt{SERVER\_\allowbreak{}ERROR}.
      \item El SJM muestra la \texttt{GUI\_\allowbreak{}MessageError} con el identificador. El bloqueo no se aplicó.
      \item Fin del caso de uso.
    \end{cuenum}} \\
   & \cuCelda{\textbf{\mbox{EX-02}: Se agota el tiempo de espera de la respuesta}\par
    \begin{cuenum}[start=1]
      \item El SJM no reenvía la solicitud y recarga la lista de bloqueados, que dirá si se aplicó.
      \item Fin del caso de uso.
    \end{cuenum}} \\
   & \cuCelda{\textbf{\mbox{EX-03}: La conexión se interrumpe durante el intercambio}\par
    \begin{cuenum}[start=1]
      \item El SJM muestra la barra de conexión perdida y recarga la lista al recuperarla.
      \item Fin del caso de uso.
    \end{cuenum}} \\
   & \cuCelda{\textbf{\mbox{EX-04}: El mensaje recibido está malformado}\par
    \begin{cuenum}[start=1]
      \item El SJM descarta el mensaje, cierra la conexión y registra en la bitácora local el contenido truncado.
      \item El SJM muestra la \texttt{GUI\_\allowbreak{}MessageError} indicando que la respuesta no pudo interpretarse.
      \item Fin del caso de uso.
    \end{cuenum}} \\ \hline
  \textbf{Postcondiciones} & \cuCelda{\begin{cureglas}
      \item \textbf{\mbox{POST-01}} Si el caso termina por el FN, existe un \textit{Bloqueo} del Jugador hacia el otro jugador, y no existe entre ellos ninguna \textit{Amistad}, \textit{Solicitud} pendiente ni \textit{Invitacion} pendiente.
      \item \textbf{\mbox{POST-02}} Si el caso termina por el \mbox{FA-02}, existe un \textit{Silencio} del Jugador hacia el otro jugador, y ninguna otra relación entre ellos cambió.
      \item \textbf{\mbox{POST-03}} Si el caso termina por el \mbox{FA-07}, la fila de \textit{Silencio} o de \textit{Bloqueo} ya no existe.
      \item \textbf{\mbox{POST-04}} El jugador bloqueado o silenciado no recibió ninguna notificación.
      \item \textbf{\mbox{POST-05}} Ninguna partida en curso se interrumpió.
    \end{cureglas}} \\ \hline
  \textbf{Reglas de negocio} & \cuCelda{\begin{cureglas}
      \item \textbf{\mbox{RN-01}} Silenciar oculta los mensajes y emotes del otro jugador a quien silencia, en todos los canales. Es reversible, no pide confirmación y no impide coincidir en partidas.
      \item \textbf{\mbox{RN-02}} Ni el silencio ni el bloqueo se notifican al afectado, y las respuestas del sistema a este no lo delatan.
      \item \textbf{\mbox{RN-03}} Bloquear elimina la amistad, las solicitudes pendientes y las invitaciones pendientes entre ambos.
      \item \textbf{\mbox{RN-04}} Un bloqueo no interrumpe una partida en curso; impide las siguientes.
      \item \textbf{\mbox{RN-05}} Un bloqueo impide emparejar, invitarse, buscarse, enviarse solicitud, compartir sala y ver el perfil del otro (\mbox{CU-16}, \mbox{CU-17}, \mbox{CU-19}, \mbox{CU-30}, \mbox{CU-32}).
      \item \textbf{\mbox{RN-06}} Desbloquear no restaura la amistad eliminada.
      \item \textbf{\mbox{RN-07}} El bloqueo no afecta a la moderación ni a los reportes.
    \end{cureglas}} \\
   & \cuCelda{\begin{cureglas}
      \item \textbf{\mbox{RN-08}} Silencio y bloqueo se eliminan físicamente al deshacerse: describen una relación presente, no un hecho histórico.
    \end{cureglas}} \\ \hline
  \textbf{Observación} & \cuCelda{\textit{Sobre por qué el silencio se guarda en el servidor.}\par
    \texttt{descripcion.md} llama local al silencio, en el sentido de que solo afecta a quien silencia. Podría haberse guardado solo en el cliente, pero entonces se perdería al cambiar de dispositivo —cosa frecuente con las sesiones simultáneas de \mbox{D-01}— y el servidor tendría que entregar mensajes que el cliente va a descartar. Guardarlo en \textit{Silencio} conserva su carácter local, porque solo lo consulta quien silenció, y lo hace sobrevivir al dispositivo.} \\ \hline
  \textbf{Datos que toca} & \cuCelda{\begin{cudatos}
      \item \textbf{\textit{Sesion}, \textit{Usuario}:} lee \textit{(paso 5)}
      \item \textbf{\textit{Bloqueo}:} lee; crea; elimina \textit{(pasos 6, 7, \mbox{FA-07})}
      \item \textbf{\textit{Silencio}:} crea; elimina \textit{(pasos \mbox{FA-02}, \mbox{FA-07})}
      \item \textbf{\textit{Amistad}, \textit{Solicitud}:} elimina las del par \textit{(paso 8)}
      \item \textbf{\textit{Invitacion}:} marca como \texttt{EXPIRADA} las pendientes del par \textit{(paso 8)}
      \item \textbf{\textit{ColaEmparejamiento}, \textit{SalaParticipante}:} lee; elimina la del retirado de la sala \textit{(paso 9)}
    \end{cudatos}} \\ \hline
\end{fichacu}
\clearpage
